\section{The ascending-torus theorem}
\label{sec:torus}

The finite diagrams constructed above will supply all the higher
finiteness properties at once.  We now state the theorem that makes
this possible and give its proof, including the passage from homology
to a classifying space with finite skeleta.  This ascending-torus
criterion and its homological construction are shared with the
companion~\cite[Theorem~2.2 and Sections~2.4--2.5]{OpenAIOvergroups}.

For a group $U$ and an endomorphism $f:U\to U$, define
\[
 T(U,f)=\langle U,t\mid t^{-1}ut=f(u)\text{ for all }u\in U\rangle.
\]
This presentation includes all the relations of $U$.  When $f$ is
injective, it is the ascending HNN extension associated to $f$.

\Needspace{18\baselineskip}
\begin{theorem}[Ascending-torus theorem]
\label{thm:torus}
Let $U$ be a group and $f:U\to U$ an injective endomorphism.  Suppose
that the following conditions hold.
\begin{enumerate}[label=\textup{(\roman*)}]
\item There are a finitely presented group $P$ and homomorphisms
      $a:U\to P$ and $b:P\to U$ such that $f=ba$.
\item For each $\ell\in\{2,3\}$ there are homomorphisms
      $h_{\ell,E}:U\to U$, indexed by the finite set
      $E\in\mathcal P_\ell$ defined above, with $h_{\ell,I}=f$.
      Whenever $E,E',E+E'\in\mathcal P_\ell$ and
      $EE'=E'E=0$, the subgroups $h_{\ell,E}(U)$ and
      $h_{\ell,E'}(U)$ commute elementwise, and
      \[
       h_{\ell,E+E'}(u)
          =h_{\ell,E}(u)h_{\ell,E'}(u)
          \qquad(u\in U).
      \]
\end{enumerate}
Then $U$ embeds naturally in $T(U,f)$, and $T(U,f)$ is of type
$\Finfty$.
\end{theorem}

The proof has three parts.  First, factoring through a finitely
presented group gives a finite presentation of the ascending torus.
Second, the two idempotent diagrams give uniformly controlled
fillings of ordinary integral cycles after applying an iterate of
$f$.  Finally, these fillings imply vanishing of homology with
products of regular coefficient modules.  A cellular argument then
constructs one classifying space with finitely many cells in every
dimension.  No word-problem hypothesis is used in any of these steps.

\begin{lemma}\label{lem:torus-presentation}
Suppose that $f=ba$ for homomorphisms $a:U\to P$ and $b:P\to U$,
where $P$ is finitely presented.  Then $T(U,f)$ is finitely presented.
This assertion does not require $f$ to be injective.
\end{lemma}

\begin{proof}
Put $g=ab:P\to P$, and write $s$ for the stable letter of $T(P,g)$.
Define maps on the respective generating groups and stable letters by
\[
 \begin{aligned}
 \Phi:T(U,f)&\longrightarrow T(P,g),
 &u&\longmapsto a(u),& t&\longmapsto s,\\
 \Psi:T(P,g)&\longrightarrow T(U,f),
 &p&\longmapsto t b(p)t^{-1},& s&\longmapsto t.
 \end{aligned}
\]
For the first map, $af=ga$ verifies the defining conjugation
relations.  For the second, the image of $s^{-1}ps$ is $b(p)$,
whereas the image of $g(p)$ is
\[
 t b(g(p))t^{-1}=t f(b(p))t^{-1}=b(p).
\]
Thus both are homomorphisms.  Their composites fix the stable letters
and satisfy
\[
 \Psi\Phi(u)=t f(u)t^{-1}=u,
 \qquad
 \Phi\Psi(p)=s g(p)s^{-1}=p.
\]
Consequently $T(U,f)\cong T(P,g)$.  If
$P=\langle y_1,\ldots,y_r\mid R_1,\ldots,R_m\rangle$, choose a word
$w_i$ representing $g(y_i)$ for each $i$.  The finite presentation
\[
 \langle y_1,\ldots,y_r,s\mid
 R_1,\ldots,R_m,\ s^{-1}y_i s=w_i\ (1\le i\le r)\rangle
\]
presents $T(P,g)$.  No injectivity of $g$ is required.
\end{proof}

The remaining task is therefore higher finiteness.  The next section
proves the uniform filling assertion; Section~\ref{sec:products}
uses it to complete the proof of Theorem~\ref{thm:torus}.
